\section{Discussion}
\label{sec:limitations}
TSD makes the frozen model's response to a reference an explicit objective
for the image latent. The reference signal and reverse sampling share one
model, and the intervention schedule provides different balances of style
strength, subject alignment, and layout retention. These properties make
TSD applicable to both composition-preserving semantic editing and freer
style-conditioned generation through a common operator.

The scope of the transfer follows the model's native reference
representation. A reference can convey object identity and layout as well
as material and texture, so the desired operating point depends on the
creative task. The current collection contains synthetic references, and
extending the evaluation to photographed materials and human-made artworks
would broaden the assessed range. Repeated teacher evaluations also make
inference cost an important design consideration. Prediction reuse and
adaptive intervention schedules are promising directions for reducing it.

\section{Conclusion}
\label{sec:conclusion}
Training-free style transfer on a frozen transformer serves Semantic Style
Editing and Global Style Transfer through a common style-locality framework.
TSD applies the native reference-conditioning score direction through
RMS-normalized detached targets for latent optimization. Interleaving those
updates with the frozen reverse trajectory provides reverse-process
regularization, without training a generator or auxiliary fake-score
network. The initial 40-style study and expanded comparisons show stronger
reference matching than native editing and direct reference guidance.
Sampling and intervention settings provide different balances of medium,
subject, and composition, with no mask or segmentation.
